\documentclass{article}
\usepackage[utf8]{inputenc}
\usepackage[T1]{fontenc}
\usepackage[spanish]{babel}
\usepackage{graphicx}
\usepackage{float}
\usepackage{array}
\usepackage{booktabs}
\usepackage{geometry}
\usepackage{fancyhdr}
\usepackage{setspace}
\usepackage{parskip}
\usepackage{amsmath}
\usepackage{amsthm} 
\usepackage{amsfonts}
\usepackage{amssymb}
\usepackage{xcolor}
\usepackage{listings}
\usepackage{caption} 
\usepackage{subcaption}
\usepackage{caption}
\usepackage{hyperref}
\usepackage{geometry}
\geometry{margin=1in}
\usepackage{tikz}
\usepackage{eso-pic}
%\usepackage{siunitx}
\usetikzlibrary{positioning, shapes.geometric, arrows.meta}
\graphicspath{ {./Imagenes/} }
\usepackage[pages=some]{background}

\lstset{
    basicstyle=\ttfamily\scriptsize, % Bajamos de small a scriptsize por lo largo de los #pragma
    breaklines=true,                 % Activa el quiebre de línea
    breakatwhitespace=false,         % Cambia a false para que rompa incluso si no hay espacios
    columns=fullflexible,            % Ayuda a que el espaciado sea más natural
    keepspaces=true,                 % Mantiene los espacios del código original
    literate={á}{{\'a}}1 {é}{{\'e}}1 {í}{{\'i}}1 {ó}{{\'o}}1 {ú}{{\'u}}1 {ñ}{{\~n}}1 % Para que acepte acentos en comentarios
}

%\usepackage[backend=biber,style=authoryear]{biblatex}
%\usepackage{csquotes}
%\addbibresource{referencias.bib}

\definecolor{codegreen}{rgb}{0,0.6,0}
\definecolor{codegray}{rgb}{0.5,0.5,0.5}
\definecolor{codeblue}{rgb}{0,0,0.95}
\definecolor{backcolour}{rgb}{0.95,0.95,0.92}

\lstdefinestyle{mystyle}{
    backgroundcolor=\color{backcolour},   
    commentstyle=\color{codegreen},
    keywordstyle=\color{codeblue},
    numberstyle=\tiny\color{codegray},
    stringstyle=\color{codegreen},
    basicstyle=\ttfamily\footnotesize,
    breakatwhitespace=false,         
    breaklines=true,                 
    captionpos=b,                    
    keepspaces=true,                 
    numbers=left,                    
    numbersep=5pt,                  
    showspaces=false,                
    showstringspaces=false,
    showtabs=false,                  
    tabsize=2
}
\lstset{style=mystyle}


\title{C\'odigos de Hamming.}
\author{Jesús Amador, Sarai Ramos, Roberto Eduardo, Belinda Santamaria}
\date{Mayo 2026}


\begin{document}

\maketitle

\section{C\'odigo de Hamming}
\noindent
  Los c\'odigos de Hamming son una familia de c\'odigos lineales de correcci\'on de 
  errores. Los c\'odigos de Hamming pueden detectar errores de uno y dos bits o corregir 
  errores de un bit sin detectar errores no corregidos.\\

\noindent
  Actualmente, los c\'odigos de Hamming son fundamentales en la teoria de la codificaci\'on y tiene una gran cantidad de aplicaciones pr\'acticas. En concreto, los c\'odigos de 
  errores tienen un papel esencial en la vida contidiana y son usados por modems, memorias e incluso en comunicaiones v\'ia sat\'elite.

\section{Distancia de Hamming}
\noindent
  Si $x,y \in \{0, 1\}^{n} = \{0, 1\}\times\{0, 1\}\times\cdots\times\{0, 1\}$, enotnces
  dis$(x,y)$ denota la distancia de Hamming entre $x$ e $y$ (es decir, la cantidad de 
  simbolos distintos). Un  c\'odigo de longitud $n$ y distancia m\'inima $d$ es un 
  subconjunto $C \subseteq \{0, 1\}^{n}$ tal que dist$(x,y) \geq d$ para todo $x,y \in C$.
  Denotemos por $A(n, d)$ el n\'umero m\'aximo de $n$-tuplas en un c\'odigo de longitud $n$ 
  con distancia m\'inima $d$.

\clearpage
%\nocite{*}
%\printbibliography[title={Referencias Bibliográficas}]
%\bibliographystyle{plain}
%\bibliography{referencias}
%\printbibliography
\end{document}



